\section{Related Work}
\label{sec:related-work}

\subsection{Autoregressive Language Models}
Autoregressive language models assign probabilities by predicting each token
from its prefix. Transformer decoders made this objective practical at scale by
combining causal self-attention with parallel training over token positions
\cite{vaswani2017attention}. The same next-token objective is still the usual
interface for large language models: provide a prompt, generate one token after
another, and sample several completions when the task may have several plausible
reasoning paths \cite{brown2020language}. Qwen2.5 and Llama~3.1, the AR
families evaluated here, follow this paradigm at modern scale
\cite{qwen2024qwen25,grattafiori2024llama}. Its appeal is practical as well as
statistical: it is straightforward to train, fits open-ended text continuation,
has well-studied scaling behavior, and is supported by mature serving systems
such as PagedAttention/vLLM \cite{kaplan2020scaling,hoffmann2022training,kwon2023pagedattention}.

The drawback is tied to the same factorization. At inference time, the model
writes sequentially; latency grows with output length, and every decision is
conditioned on a prefix that has already been committed. Sampling several
chains is one way to soften that commitment. Chain-of-thought prompting and
self-consistency show that repeated samples can improve reasoning by pooling
answers across different solution traces \cite{wei2022chain,wang2023selfconsistency}.
Pass@k, first popularized in code-generation evaluation, captures the same
idea as the probability that at least one of $k$ samples is correct
\cite{chen2021humaneval}. Large-$k$ reasoning evaluations also reveal ranking
changes that pass@1 hides: Yue et al. report that RLVR-trained models can lead
at small $k$, while their base models can overtake at large $k$ because the
sampled sets cover more solutions \cite{yue2025rlvr}. For this thesis, that
result is the key methodological warning: a single sampled answer is not enough
to characterize how a reasoning system behaves.

\subsection{Non-Autoregressive, Masked, and Iterative Text Generation}
Diffusion language models build on a longer line of bidirectional and
iterative text generation. Masked language modeling trains a model to recover
hidden tokens using context from both sides; BERT is not a left-to-right
generator, but it made this bidirectional view central to modern NLP
\cite{devlin2019bert}. Conditional masked language models then adapted the idea
for generation by repeatedly replacing low-confidence tokens, so several
positions could be updated in parallel instead of being fixed strictly from
left to right \cite{ghazvininejad2019maskpredict}. That line of work already
sets up the main design contrast used here: generation as sequential commitment
versus generation as iterative refinement.

The masked and iterative view is especially relevant for structured outputs. Because later refinement steps can condition on partial content on both sides, non-autoregressive methods are naturally suited to infilling, editing, and constraint satisfaction. However, early masked-generation systems were typically evaluated on narrower conditional-generation tasks and did not match the scale, instruction-following ability, or infrastructure maturity of modern autoregressive LLMs. Recent diffusion-based LLMs can therefore be understood not as an isolated departure, but as a large-scale continuation of masked denoising and iterative refinement.

\subsection{Diffusion-Based Language Models}
Diffusion models generate by reversing a corruption process. In text, this idea is implemented either through continuous relaxations of token embeddings or through discrete corruption processes over categorical states. D3PMs showed how structured transition matrices and absorbing mask states can extend denoising diffusion to discrete data, including text, and connect diffusion to mask-based generation \cite{austin2021structured}. Diffusion-LM demonstrated that continuous diffusion over word embeddings can support fine-grained controllable text generation through iterative denoising \cite{li2022diffusionlm}. DiffusionBERT linked the same family of ideas more directly to generative masked language modeling by training a BERT-like model to reverse a masked diffusion process \cite{he2023diffusionbert}.

Large diffusion language models scale this denoising perspective to contemporary LLM settings. LLaDA trains a large masked diffusion model from scratch with pre-training and supervised fine-tuning, showing that instruction following and benchmark performance are not inherently tied to autoregressive decoding \cite{nie2025llada}. A complementary line adapts pretrained autoregressive checkpoints into diffusion models, reducing the cost of training diffusion LMs at scale and clarifying objective-level connections between the paradigms \cite{gong2025scaling}. Dream follows this practical route: it is an open 7B diffusion LLM initialized from an autoregressive model and reports strong performance on general, mathematical, coding, and planning-style tasks, with inference modes that include arbitrary-order generation, infilling, and tunable quality--speed trade-offs \cite{ye2025dream}. Recent surveys group these developments around a recurring set of proposed advantages: bidirectional context, iterative refinement, parallel or semi-parallel decoding, controllability, and editing flexibility \cite{li2025survey,yu2025discrete}. These properties motivate evaluation, but they do not by themselves imply better reasoning or planning performance.

These advantages should be read as conditional rather than automatic. The quality and efficiency of a diffusion LM depend on the noise schedule, number of denoising steps, block length, confidence thresholding, cache strategy, model scale, and post-training procedure. In particular, the practical latency benefit of diffusion depends on whether parallel token updates compensate for repeated denoising passes and whether the inference system can reuse computation effectively.

\subsection{Comparative Studies of Autoregressive and Diffusion Language Models}
Recent comparisons have moved beyond the basic question of whether diffusion
LMs can produce fluent text. Work that adapts autoregressive models into
diffusion models reports competitive language modeling, reasoning, commonsense,
infilling, and instruction-following results, while also showing why clean
comparison is hard: training data, initialization, and objectives may all
differ across paradigms \cite{gong2025scaling}. LLaDA and Dream make a similar
point from the model side. They show that diffusion LMs can be competitive on
broad benchmark suites, but they do not by themselves isolate whether a gain
comes from denoising, scale, data, post-training, or the benchmark interface
\cite{nie2025llada,ye2025dream}.

Other work examines dimensions that are easy to miss in a single benchmark
table. A text-embedding comparison argues that bidirectional diffusion
architectures can be better aligned with retrieval and long-document
representation than the unidirectional attention inherited from autoregressive
pre-training \cite{zhang2025embedding}. Systems-oriented analyses study latency
and throughput, finding that diffusion can use more sequence-level parallelism
but may struggle with long contexts or batched serving unless decoding schedules
and caching are designed carefully \cite{kim2025beyond}. Fast-dLLM addresses
part of this deployment gap through approximate KV caching and confidence-aware
parallel decoding \cite{wu2025fastdllm}. Taken together, this literature makes
the comparison multi-dimensional: generation order, objective, decoding budget,
context length, controllability, latency, throughput, and evaluation protocol
can each change the conclusion.

\subsection{Evaluation Gaps and Motivation for This Work}
Prior work shows that diffusion language models can generate fluent text and
can be competitive on several benchmark suites. What remains less clear is how
to compare them with autoregressive models on reasoning and planning tasks.
Existing studies often vary model size, tokenizer, initialization, training
data, prompt format, denoising schedule, sampling budget, or inference backend
at the same time. When those factors move together, it is hard to tell whether
an observed gain comes from the generation interface or from the surrounding
experimental setup.

Reasoning benchmarks complicate the picture further. A single deterministic
answer and a many-sample pass@k score reward different things. The first asks
whether one decoded solution is correct; the second asks whether the candidate
set contains at least one correct answer. Yue et al.'s RLVR/base-model reversal
is a concrete example: a model that looks weaker under top-1 evaluation can
still reveal more correct candidates when the sampling budget is large
\cite{yue2025rlvr}. Parser behavior can also shift measured accuracy,
especially when a task expects a specific output format. Aggregate scores then
hide another question: whether different models solve the same instances or
different ones.

These gaps shape the evaluation design in this thesis. Before interpreting a
possible paradigm difference, the comparison needs to make the prompt interface,
denoising or sampling budget, parser behavior, and benchmark structure explicit.
The goal is therefore not only to ask which model has the highest score, but to
track how that score changes when the evaluator can inspect more sampled
candidates.
